% الفصل 22 — الزمن الحقيقيّ بعمق: WebSocket و SignalR وفكرة الـ Hubs
\section{الزمن الحقيقيّ بعمق: \en{WebSocket} و\en{SignalR} وفكرة الـ \en{Hubs}}

في الجزء الأول رأينا الإشعارات والشات وهما يشتغلان: رسالة تظهر عند الطرف
الآخر لحظيًّا، وجرس الإشعارات يرنّ بلا تحديث للصفحة. هذا الفصل يفتح الغطاء ويجيب
عن السؤال الأصعب: كيف يصل التحديث من الخادم إلى المتصفّح في نفس اللحظة،
بلا أن يسأل المتصفّح أصلًا؟ هنمشي «واحدة واحدة»: نبدأ من مشكلة الـ \en{HTTP}
العاديّ، ثم نفهم الـ \en{WebSocket} كحلّ، ثم نرى \en{SignalR} كطبقة تجريد فوقه،
وأخيرًا نغوص في الكود الحقيقيّ لثلاثة \en{hubs} في \en{ScholarPath}: كيف يتوثّق
الاتصال، وكيف يُوجَّه الدفع بالمجموعات، وأين يجلس «المشتركون» على الطرفين.

هذه المفاهيم من المفترض أنّك رأيتها مستخدَمة في الجزء الأول؛ هنا نشرح الآليّة
تحتها لا الاستخدام.

% ─────────────────────────────────────────────────────────────
\subsection{المشكلة أولًا: لماذا لا يكفي \en{HTTP} العاديّ؟}

الـ \en{HTTP} التقليديّ يعمل بنمط «طلب ثم ردّ»: المتصفّح يفتح اتصالًا، يرسل
طلبًا، يستقبل الردّ، ثم يُغلَق الاتصال. الخادم لا يملك أي طريقة ليبدأ
الحديث من نفسه؛ هو يردّ فقط حين يُسأل. طيب، ماذا لو وصلت رسالة شات جديدة وأنا لا
أسأل؟ لن أعرف بها.

الحلّ الساذج اسمه \en{polling}: اجعل المتصفّح يسأل الخادم كل ثانية «في جديد؟».
هذا يعمل، لكنّه مكلف: آلاف الطلبات الفارغة (\en{no data}) لكل مستخدم، كل واحد منها
يحمل ترويسات \en{HTTP} كاملة ومصافحة \en{TLS}، وتأخير يصل إلى ثانية كاملة قبل أن
يرى المستخدم أي تحديث. تخيّلها مضروبة في آلاف المتّصلين.

\begin{why}
اللحظيّة الحقيقيّة تحتاج أن يدفع الخادم البيانات متى شاء — لا أن ينتظر
سؤالًا. الـ \en{polling} يقلب المعادلة: يجعل العميل يسأل باستمرار، فيدفع
ثمن التأخير والحِمل معًا. نحتاج قناةً يبقى الخادم قادرًا فيها على الكلام.
\end{why}

% ─────────────────────────────────────────────────────────────
\subsection{\en{WebSocket}: مصافحة تُرقّي الاتصال إلى قناة دائمة}

الـ \en{WebSocket} بروتوكول يحلّ هذا بالضبط. يبدأ حياته كطلب \en{HTTP} عاديّ، لكن
بترويسة خاصة \code{Upgrade: websocket}. الخادم يوافق بردّ الحالة
\code{101 Switching Protocols}، وهذه اللحظة تُسمّى «المصافحة» \en{(handshake)}.
بعدها لا يُغلَق الاتصال، بل يُرقّى إلى قناة دائمة ثنائية الاتجاه
\en{(full-duplex)}: الطرفان يرسلان في أي وقت وعبر نفس الاتصال المفتوح.

النتيجة المهمة: صار الخادم قادرًا على أن يدفع \en{(push)} رسالةً للعميل بلا
أي طلب سابق. لا مزيد من \en{polling}، ولا تأخير، واتصال واحد يخدم كل التحديثات.

\begin{center}\begin{LTR}
\begin{tikzpicture}
  % ---- polling (المشكلة) ----
  \node[actor] (cp) at (0,0) {Client};
  \node[actor] (sp) at (6,0) {Server};
  \draw[arr] ([yshift=5pt]cp.east) -- node[lbl,above]{GET /updates?} ([yshift=5pt]sp.west);
  \draw[arrd] ([yshift=-7pt]sp.west) -- node[lbl,below]{204 no data} ([yshift=-7pt]cp.east);
  \node[fbox,text width=6cm] at (3,-1.25) {polling: يتكرّر كل ثانية = حِمل + تأخير};

  % ---- websocket (الحل) ----
  \node[actor] (cw) at (0,-3.1) {Client};
  \node[actor] (sw) at (6,-3.1) {Server};
  \draw[arr] ([yshift=9pt]cw.east) -- node[lbl,above]{HTTP Upgrade (handshake)} ([yshift=9pt]sw.west);
  \draw[arrb] (cw.east) -- (sw.east);
  \draw[arr] ([yshift=-9pt]sw.west) -- node[lbl,below]{server push (بلا طلب)} ([yshift=-9pt]cw.east);
  \node[tbox,text width=6cm] at (3,-4.4) {قناة واحدة دائمة ثنائية الاتجاه (full-duplex)};
\end{tikzpicture}
\end{LTR}\end{center}

في الأعلى: الـ \en{polling} يرسل طلبًا بعد طلب فيرجع أغلبها فارغًا. في الأسفل:
مصافحة واحدة تفتح قناة دائمة، ثم يدفع الخادم متى شاء عبرها.

% ─────────────────────────────────────────────────────────────
\subsection{\en{SignalR} فوق الـ \en{WebSocket}: تجريد يخفي النقل}

الـ \en{WebSocket} قويّ لكنّه منخفض المستوى: أنت تدير أُطُر البايتات
\en{(frames)}، وإعادة الاتصال عند القطع، وماذا تفعل على المتصفّحات أو الشبكات
التي تحظر الـ \en{WebSocket}. هنا يدخل \en{SignalR}: مكتبة من \en{ASP.NET Core}
تجلس فوق الـ \en{WebSocket} وتخفي كل هذا التعقيد.

\begin{idea}[الفكرة: \en{SignalR} = تجريد النقل + إدارة الاتصالات]
\en{SignalR} يفاوض تلقائيًّا \en{(negotiate)} على أفضل نقل متاح: يستخدم الـ
\en{WebSocket} أولًا، وإن تعذّر يسقط \en{(fallback)} إلى \en{Server-Sent Events}
ثم \en{Long Polling}. كودك لا يتغيّر إطلاقًا مهما كان النقل تحته. زيادةً على ذلك،
يدير \en{SignalR} تعريف كل اتصال \en{(ConnectionId)}، وتجميع الاتصالات في
مجموعات، وإعادة الاتصال بعد القطع.
\end{idea}

القلب المفاهيميّ في \en{SignalR} اسمه الـ \en{Hub}: صنف \en{C\#} على الخادم يعمل
كنقطة نداء ثنائية الاتجاه. العميل يقدر يستدعي دوالّ هذا الصنف (مثل
\code{JoinConversation})، والخادم يقدر يدفع أحداثًا مسمّاة للعميل (مثل
\code{"notification"}). كأنّه هاتف: كلا الطرفين يرفع السماعة ويتكلّم.

\begin{keyp}[نقطة مهمة: الـ \en{Hub} ثنائيّ الاتجاه]
لا تتخيّل الـ \en{Hub} كـ \en{controller} يردّ على الطلبات فقط. هو قناة في
الاتجاهين: صعودًا العميل ينادي دوالّ الـ \en{Hub}، ونزولًا
الخادم يدفع أحداثًا عبر \code{Clients} إلى عميل أو مجموعة عملاء.
\end{keyp}

في \en{ScholarPath} ثلاثة \en{hubs} مربوطة على مسارات ثابتة: \code{/hubs/chat}
و\code{/hubs/notifications} و\code{/hubs/community}.

% ─────────────────────────────────────────────────────────────
\subsection{دورة حياة الاتصال والهُويّة من الـ \en{JWT}}

لكل \en{hub} دالّتان تُستدعيان تلقائيًّا: \code{OnConnectedAsync} لحظة اتصال
عميل جديد، و\code{OnDisconnectedAsync} لحظة انقطاعه. هذا المكان المثاليّ لضمّ
الاتصال إلى مجموعاته أو إعلان حالة الحضور.

لكن قبل أي شيء: مَن هذا المتّصل؟ كل \en{hubs} النظام ترث من صنف أساس واحد يفرض
التوثيق. هنشوف الملف الحقيقيّ.

\begin{english}
\begin{lstlisting}[language={[Sharp]C}]
[Authorize]                              // only authenticated connections
public abstract class AuthenticatedHub : Hub
{
    public Task<string> Ping() =>
        Task.FromResult($"pong from {GetType().Name}");
}

public sealed class NotificationHub : AuthenticatedHub
{
    public override async Task OnConnectedAsync()
    {
        // AddToGroup MUST be awaited (see the last section on the race)
        if (Context.UserIdentifier is not null)
        {
            await Groups.AddToGroupAsync(
                Context.ConnectionId, $"user:{Context.UserIdentifier}");
        }
        await base.OnConnectedAsync();
    }
}
\end{lstlisting}
\end{english}

الشرح سطرًا بسطر لأهم الأسطر:
\begin{itemize}
\item \code{[Authorize]} فوق \code{AuthenticatedHub}: بوّابة توثيق على مستوى
  الاتصال كلّه. أي عميل يحاول فتح قناة بلا \en{JWT} صالح يُرفَض في المصافحة نفسها،
  فلا يصل إلى أي دالّة.
\item \code{Context.UserIdentifier}: هوية المتّصل مستخرَجة تلقائيًّا من الـ
  \en{JWT} (من مطالبة المعرّف). هذه القيمة موثوقة لأن الـ \en{token} موقَّع؛
  العميل لا يستطيع تزويرها.
\item \code{Groups.AddToGroupAsync(Context.ConnectionId, "user:\{id\}")}: بمجرّد
  الاتصال، نضمّ هذا الاتصال إلى مجموعته الشخصية \code{user:\{id\}}. بعد هذا السطر
  يصبح العميل جاهزًا لاستقبال أي إشعار يُدفَع إلى تلك المجموعة.
\end{itemize}

\begin{secnote}
التوثيق على مستوى الـ \en{hub} فارق أمنيّ جوهريّ: القناة الحيّة ليست ثغرة خلفية
تتجاوز التوثيق. نفس الـ \en{JWT} الذي يحمي الـ \en{REST API} يحمي الـ
\en{WebSocket}، والهُويّة تأتي من الـ \en{token} لا من أي معطى يرسله العميل.
\end{secnote}

الرسم التالي يلخّص دورة الحياة من الاتصال إلى الجاهزية للاستقبال:

\begin{center}\begin{LTR}
\begin{tikzpicture}[node distance=0.42cm]
  \node[wbox] (s1) {Client connects: withUrl("/hubs/notifications") + JWT};
  \node[wbox,below=0.42cm of s1] (s2) {[Authorize] gate: verify JWT, set Context.UserIdentifier};
  \node[wbox,below=0.42cm of s2] (s3) {OnConnectedAsync()};
  \node[wbox,below=0.42cm of s3] (s4) {await Groups.AddToGroupAsync(ConnectionId, "user:id")};
  \node[wbox,below=0.42cm of s4,fill=spLightTeal] (s5) {جاهز لاستقبال أي push فوريًّا};
  \draw[arr] (s1)--(s2); \draw[arr] (s2)--(s3);
  \draw[arr] (s3)--(s4); \draw[arr] (s4)--(s5);
\end{tikzpicture}
\end{LTR}\end{center}

% ─────────────────────────────────────────────────────────────
\subsection{المجموعات والتوجيه \en{(Groups \& Routing)}}

المجموعة \en{(Group)} في \en{SignalR} مجرّد اسم نصّي يجمع تحته عدّة اتصالات.
الدفع يصير للاسم لا للاتصال الفرديّ، وهذا يحلّ مشكلتين: المستخدم الواحد قد يفتح
عدّة تبويبات (عدّة اتصالات، مجموعة واحدة \code{user:\{id\}})، والحدث الواحد قد
يخصّ عدّة مستخدمين (مثل محادثة أو تصنيف منتدى).

\en{ScholarPath} يستعمل ثلاث عائلات أسماء واضحة:

\begin{center}
\begin{tabular}{L{4cm} L{8cm}}
\toprule
{\color{spPrimary}المجموعة} & {\color{spPrimary}المعنى} \\
\midrule
\code{user:\{id\}} & كل اتصالات مستخدم واحد — قناته الشخصية للإشعارات \\
\code{conversation:\{id\}} & طرفا محادثة شات — رسائل ومؤشّر الكتابة \\
\code{forum-category:\{slug\}} & متابعو تصنيف في المنتدى — مواضيع جديدة \\
\code{forum-post:\{id\}} & مشاهدو موضوع واحد — الردود الجديدة \\
\bottomrule
\end{tabular}
\end{center}

أما التوجيه فيتمّ عبر خصائص \code{Clients} داخل الـ \en{hub}:
\code{Clients.Group(name)} للدفع لمجموعة كاملة، و\code{Clients.User(id)} لمستخدم
بعينه، و\code{Clients.Others} لكل المتّصلين ما عدا المُرسِل، و
\code{Clients.OthersInGroup(name)} لكل أعضاء المجموعة ما عدا المُرسِل — مثاليّ
لمؤشّر «يكتب الآن» حتى لا يرى المرسِل نفسه يكتب.

خلّينا نشوف \code{ChatHub} حقيقيًّا: كيف يضمّ اتصالًا لمحادثة، وكيف يبثّ الكتابة،
والأهمّ: بوّابة تمنع التنصّت.

\begin{english}
\begin{lstlisting}[language={[Sharp]C}]
public async Task JoinConversation(string conversationId)
{
    await EnsureParticipantAsync(conversationId);          // IDOR gate
    await Groups.AddToGroupAsync(
        Context.ConnectionId, $"conversation:{conversationId}");
}

public async Task TypingStart(string conversationId)
{
    await EnsureParticipantAsync(conversationId);
    await Clients.OthersInGroup($"conversation:{conversationId}")
        .SendAsync("TypingStart", conversationId, Context.UserIdentifier);
}

private async Task EnsureParticipantAsync(string conversationId)
{
    // the JWT caller must be one of the two participants
    // ... parse ids ...
    var isParticipant = await db.Conversations.AsNoTracking()
        .AnyAsync(c => c.Id == convId
            && (c.ParticipantOneId == userId || c.ParticipantTwoId == userId));
    if (!isParticipant)
        throw new HubException("You are not a participant.");
}
\end{lstlisting}
\end{english}

الشرح سطرًا بسطر لأهم الأسطر:
\begin{itemize}
\item \code{JoinConversation}: العميل ينادي هذه الدالّة صعودًا حين يفتح محادثة.
  لاحظ أنّها تستدعي \code{EnsureParticipantAsync} قبل الضمّ للمجموعة.
\item \code{Clients.OthersInGroup(...)}: مؤشّر الكتابة يُبثّ لكل أعضاء
  \code{conversation:\{id\}} ما عدا المُرسِل — فلا أرى نفسي «أكتب».
\item \code{EnsureParticipantAsync}: البوّابة الأمنية. تتأكّد من قاعدة البيانات
  أنّ صاحب الـ \en{JWT} أحد طرفَي المحادثة فعلًا؛ وإلا ترمي \code{HubException}.
\end{itemize}

\begin{spotlight}[ركّز هنا لإبهار اللجنة: بوّابة \en{IDOR} على القناة الحيّة]
اسم المجموعة \code{conversation:\{id\}} قابل للتخمين. لولا
\code{EnsureParticipantAsync} لاستطاع أي مستخدم موثَّق أن ينادي
\code{JoinConversation} بمعرّف محادثة غيره ويتنصّت على رسائل خاصّة تُبثّ لتلك
المجموعة. الدرس: التوجيه بالمجموعات لا يكفي وحده كأمان — لا بدّ من بوّابة
تتحقّق أنّك طرفٌ فعليّ قبل الانضمام. نفس المنطق مطبَّق في \code{CommunityHub} على
التصنيفات والمواضيع (تصنيف نشِط، وموضوع مرئيّ غير محذوف).
\end{spotlight}

% ─────────────────────────────────────────────────────────────
\subsection{أين المشتركون؟ على الخادم عبر \en{IHubContext}، وعلى العميل عبر \code{.on}}

سؤال جوهريّ: الدفع رأيناه داخل الـ \en{hub}. لكن الإشعار يُولَد عادةً في
مكان بعيد — بعد نجاح دفعة، أو قبول حجز. كيف ندفع من خارج الـ \en{hub}؟ الجواب:
\code{IHubContext<THub>}، وهو مقبض للـ \en{hub} تحقنه أي خدمة عبر حقن التبعيّات.

هذا بالضبط ما يفعله \code{NotificationDispatcher}: يحفظ صفوف الإشعار في قاعدة
البيانات، ثم يدفع النسخة الحيّة عبر \code{IHubContext} إلى مجموعة المستخدم.

\begin{english}
\begin{lstlisting}[language={[Sharp]C}]
public sealed class NotificationDispatcher(
    IHubContext<NotificationHub> hub, /* ... */) : INotificationDispatcher
{
    // ... persist a Notification row per channel first ...

    private async Task DeliverAsync(Notification row, /* ... */)
    {
        // InApp: real-time toast to the recipient's personal group
        if (!silenced)
            await hub.Clients.Group($"user:{row.RecipientUserId}")
                .SendAsync("notification", new
                {
                    id = row.Id,
                    type = row.Type.ToString(),
                    row.TitleAr, row.BodyAr, row.DeepLink,
                }, ct);
    }
}
\end{lstlisting}
\end{english}

الشرح: \code{IHubContext<NotificationHub>} محقون في الباني \en{(constructor)}؛
الخدمة لا ترث من الـ \en{hub} ولا تعرف تفاصيله، فقط تدفع عبره. السطر
\code{hub.Clients.Group("user:\{id\}").SendAsync("notification", ...)} يبثّ حدثًا
اسمه \code{"notification"} لمجموعة المستخدم، فيصل لكل تبويباته المفتوحة دفعةً
واحدة. لاحظ الحقول العربية \code{TitleAr} و\code{BodyAr}: الحمولة ثنائية اللغة.

أما «المشترك» على الطرف الآخر، فيجلس على العميل ويربط اسم الحدث بدالّة معالجة عبر
\code{connection.on("event", handler)}. هنشوف كود الواجهة الحقيقيّ (\en{TypeScript}):

\begin{english}
\begin{lstlisting}[language=Java]
const connection = new signalR.HubConnectionBuilder()
  .withUrl(apiBase + "/hubs/notifications", {
    // current token on every (re)connect — survives token refresh
    accessTokenFactory: () =>
      useAuthStore.getState().tokens?.accessToken ?? "",
    // offer all transports so it works where WebSockets are blocked
    transport:
      signalR.HttpTransportType.WebSockets |
      signalR.HttpTransportType.ServerSentEvents |
      signalR.HttpTransportType.LongPolling,
  })
  .withAutomaticReconnect()
  .build();

// subscriber: fires whenever the server pushes "notification"
connection.on("notification", (payload) => {
  toast(payload.titleAr, { description: payload.bodyAr });
});
await connection.start();
\end{lstlisting}
\end{english}

الشرح سطرًا بسطر لأهم الأسطر:
\begin{itemize}
\item \code{withUrl(apiBase + "/hubs/notifications")}: نفس مسار الـ \en{hub}
  المربوط على الخادم؛ الـ \en{apiBase} يتبع الـ \en{API} المنشور على \en{Azure}.
\item \code{accessTokenFactory}: يُقرأ الـ \en{token} الحاليّ من المخزن في كل
  اتصال أو إعادة اتصال، فلا ينكسر الاتصال بعد تدوير الـ \en{token}.
\item \code{transport: WebSockets | ServerSentEvents | LongPolling}: نعرض كل
  وسائل النقل ونُبقي خطوة التفاوض، فلو تعذّر الـ \en{WebSocket} يسقط تلقائيًّا.
\item \code{connection.on("notification", ...)}: هنا المشترك. أي دفعة يبثّها
  \code{NotificationDispatcher} بهذا الاسم تُشغّل هذه الدالّة فورًا فيظهر الـ
  \en{toast}.
\end{itemize}

الرسم التالي يجمع الصورة: حدث واحد على الخادم يتفرّع \en{(fan-out)} إلى كل عملاء
المجموعة.

\begin{center}\begin{LTR}
\begin{tikzpicture}[node distance=0.7cm and 1.25cm]
  \node[srcbox,text width=2.7cm] (evt) {NotificationDispatcher\\(server event)};
  \node[tbox,right=1.25cm of evt,text width=2.7cm] (ctx) {IHubContext\\Clients.Group(...)};
  \node[gbox,right=1.25cm of ctx,text width=2.1cm] (grp) {group\\user:id};
  \node[actor,above right=0.35cm and 1.25cm of grp] (c1) {browser tab};
  \node[actor,right=1.25cm of grp] (c2) {second tab};
  \node[actor,below right=0.35cm and 1.25cm of grp] (c3) {phone};
  \draw[arr] (evt) -- node[lbl,above]{SendAsync} (ctx);
  \draw[arr] (ctx) -- (grp);
  \draw[arr] (grp) -- (c1);
  \draw[arr] (grp) -- (c2);
  \draw[arr] (grp) -- (c3);
\end{tikzpicture}
\end{LTR}\end{center}

\begin{spotlight}[ركّز هنا لإبهار اللجنة: اتصال واحد دائم يخدم كل التحديثات]
كل مستخدم يفتح اتصالًا حيًّا واحدًا لكل \en{hub}، ويبقى مفتوحًا طوال
الجلسة. الإشعارات، والشات، ومؤشّر الكتابة، والحضور، وردود المنتدى — كلها تسري في
نفس القناة الدائمة. صفر \en{polling}، صفر تأخير، وحِمل أدنى بكثير من سؤال
الخادم كل ثانية. هذا هو الفرق بين واجهة «تحسّها حيّة» وأخرى تحتاج تحديثًا يدويًّا.
\end{spotlight}

% ─────────────────────────────────────────────────────────────
\subsection{الـ \en{hubs} الثلاثة: خريطة الجروبات والأحداث المدفوعة}

نلخّص المعمار كلّه في جدول واحد: كل \en{hub}، ومساره، ومجموعاته، وأحداثه المدفوعة
نزولًا إلى العميل.

\begin{center}
\begin{tabular}{L{2.9cm} L{3.9cm} L{5.2cm}}
\toprule
{\color{spPrimary}الـ \en{Hub}} & {\color{spPrimary}المجموعات} & {\color{spPrimary}أحداث تُدفَع للعميل} \\
\midrule
\code{NotificationHub}\newline \code{/hubs/notifications} & \code{user:\{id\}} &
\code{"notification"} (إشعار حيّ لكل تبويبات المستخدم) \\
\code{ChatHub}\newline \code{/hubs/chat} & \code{conversation:\{id\}} &
\code{"MessageReceived"}، \code{"TypingStart"}، \code{"TypingStop"}،
\code{"UserOnline"}، \code{"UserOffline"} \\
\code{CommunityHub}\newline \code{/hubs/community} & \code{forum-category:\{slug\}}،
\code{forum-post:\{id\}} & \code{"ReplyCreated"} وتحديثات المواضيع الحيّة \\
\bottomrule
\end{tabular}
\end{center}

\begin{note}
لاحظ الاتساق: الأحداث النازلة أسماء نصّية (\code{"MessageReceived"}) يشترك فيها
العميل عبر \code{.on}، والدوالّ الصاعدة (\code{JoinConversation}) يناديها العميل
عبر \code{invoke}. اسم الحدث على الطرفين يجب أن يتطابق حرفيًّا؛ خطأ حرف واحد =
مشترك لا يسمع شيئًا.
\end{note}

% ─────────────────────────────────────────────────────────────
\subsection{تفصيلة دقيقة: لماذا \code{await} على \code{AddToGroupAsync}؟}

عُد لكود \code{NotificationHub.OnConnectedAsync}: الضمّ للمجموعة مسبوق بـ
\code{await}. لماذا لا نتركه «أطلِق وانسَ» \en{(fire-and-forget)}؟ لأن هذا يفتح
باب \en{race condition} خبيث.

\begin{qa}[سؤال متوقّع: ماذا يحدث لو لم نضع \code{await}؟]
تخيّل السيناريو: العميل اتصل، وبدأ ضمّه للمجموعة \code{user:\{id\}} في الخلفية بلا
انتظار. في نفس اللحظة أطلق حدثٌ آخر إشعارًا فوريًّا لتلك المجموعة. لو سبق الإشعارُ
اكتمالَ الانضمام، يُدفَع لمجموعة لا يوجد بها هذا الاتصال بعد — فيضيع
الإشعار الحيّ. المستخدم لن يراه إلا بعد إعادة اتصال لاحقة، وربما يظهر عندها كأنّه
«مكرَّر». وضع \code{await} يضمن أن الاتصال صار عضوًا كامل العضوية قبل أن
يعود التحكّم، فلا نافذة سباق.
\end{qa}

هذه التفصيلة الصغيرة درسٌ عام: في الكود غير المتزامن، ترتيب اكتمال العمليّات ليس
تفصيلًا تجميليًّا — قد يكون الفرق بين إشعار يصل وإشعار يضيع.

% ─────────────────────────────────────────────────────────────
\begin{summary}[الخلاصة في سطور]
\begin{itemize}
\item الـ \en{HTTP} العاديّ طلب-وردّ ثم يُغلَق؛ لِعرض التحديثات نحتاج \en{polling}
  مكلفًا وبطيئًا لأن الخادم لا يبدأ الحديث.
\item الـ \en{WebSocket} يبدأ بمصافحة \en{(HTTP Upgrade)} تُرقّي الاتصال إلى قناة
  دائمة ثنائية الاتجاه، فيدفع الخادم \en{(push)} متى شاء بلا سؤال.
\item \en{SignalR} تجريد فوق الـ \en{WebSocket}: يفاوض على أفضل نقل ويسقط للأقدم
  عند الحاجة، ويدير الاتصالات والمجموعات. القلب هو الـ \en{Hub}: قناة ثنائية —
  العميل ينادي دوالّه، والخادم يدفع أحداثه.
\item التوثيق عبر \code{[Authorize]} على الـ \en{hub}، والهُويّة من الـ \en{JWT}
  عبر \code{Context.UserIdentifier} (لا يمكن تزويرها).
\item التوجيه بالمجموعات: \code{user:\{id\}}، \code{conversation:\{id\}}،
  \code{forum-category:\{slug\}}؛ والدفع عبر \code{Clients.Group} / \code{User} /
  \code{Others} / \code{OthersInGroup}.
\item المشتركون: على العميل عبر \code{connection.on("event")}، وعلى الخادم من خارج
  الـ \en{hub} عبر \code{IHubContext} (كما في \code{NotificationDispatcher} و
  \code{ChatRealtimeNotifier}).
\item بوّابة \en{IDOR} (لازم تكون طرفًا) تمنع التنصّت على قنوات الغير، و
  \code{await} على \code{AddToGroupAsync} يمنع سباقًا يفوّت إشعارًا.
\end{itemize}
\end{summary}
